\documentclass[a4paper,10pt,landscape]{article}

\usepackage[utf8]{inputenc}
\usepackage[T1]{fontenc}
\usepackage[ngerman]{babel}
\usepackage[landscape, top=0.8cm, bottom=0.8cm, left=0.9cm, right=0.9cm]{geometry}
\usepackage{xcolor}
\usepackage{listings}
\usepackage{tcolorbox}
\usepackage{multicol}
\usepackage{fontawesome5}
\usepackage{tikz}
\usepackage{microtype}
\usepackage{lmodern}
\usepackage{parskip}

\setlength{\parskip}{2pt}
\setlength{\columnsep}{0.7cm}
\setlength{\columnseprule}{0.4pt}
\def\columnseprulecolor{\color{arduinoBlue!30}}

% ── Farben ──────────────────────────────────────────────────────────────
\definecolor{arduinoBlue}{HTML}{006D6F}
\definecolor{arduinoLightBlue}{HTML}{00A0A3}
\definecolor{goodGreen}{HTML}{2E7D32}
\definecolor{badRed}{HTML}{B71C1C}
\definecolor{codeBackground}{HTML}{F3F4F6}
\definecolor{headerBg}{HTML}{006D6F}
\definecolor{commentColor}{HTML}{5C6BC0}
\definecolor{keywordColor}{HTML}{1565C0}
\definecolor{stringColor}{HTML}{2E7D32}
\definecolor{ruleAccent}{HTML}{00ACC1}
\definecolor{ruleAccent2}{HTML}{00C174}
\definecolor{ruleAccent3}{HTML}{C1AA00}
\definecolor{ruleAccent4}{HTML}{C17700}

% ── Listings ────────────────────────────────────────────────────────────
\lstdefinestyle{cs}{
  language=C++,
  backgroundcolor=\color{codeBackground},
  basicstyle=\ttfamily\fontsize{7}{8.5}\selectfont,
  keywordstyle=\color{keywordColor}\bfseries,
  commentstyle=\color{commentColor}\itshape,
  stringstyle=\color{stringColor},
  numbers=none,
  frame=none,
  breaklines=true,
  showstringspaces=false,
  tabsize=2,
  xleftmargin=4pt,
  xrightmargin=2pt,
  morekeywords={void,int,bool,byte,float,const,String,
                pinMode,digitalWrite,digitalRead,
                analogWrite,analogRead,delay,
                Serial,begin,println,print,
                INPUT,OUTPUT,HIGH,LOW,setup,loop},
  literate={ä}{{\"a}}1 {ö}{{\"o}}1 {ü}{{\"u}}1
           {Ä}{{\"A}}1 {Ö}{{\"O}}1 {Ü}{{\"U}}1 {ß}{{\ss}}1
}
\lstset{style=cs}

% ── tcolorbox ───────────────────────────────────────────────────────────
\tcbuselibrary{breakable}

% Regelbox (ganzer Block)
\newtcolorbox{regelbox}[2]{
  colback=#1!6,
  colframe=#1,
  coltitle=white,
  fonttitle=\bfseries\fontsize{8}{9}\selectfont,
  title={#2},
  arc=3pt,
  boxrule=1pt,
  top=3pt, bottom=3pt, left=4pt, right=4pt,
  toptitle=2pt, bottomtitle=2pt
}

% Gut / Schlecht nebeneinander
\newtcolorbox{goodbox}{
  colback=goodGreen!7,
  colframe=goodGreen,
  fonttitle=\bfseries\fontsize{7}{8}\selectfont,
  title={\faCheck\enspace Richtig},
  arc=2pt, boxrule=0.8pt,
  top=2pt, bottom=2pt, left=3pt, right=3pt,
  toptitle=1pt, bottomtitle=1pt
}

\newtcolorbox{badbox}{
  colback=badRed!7,
  colframe=badRed,
  fonttitle=\bfseries\fontsize{7}{8}\selectfont,
  title={\faTimes\enspace Falsch},
  arc=2pt, boxrule=0.8pt,
  top=2pt, bottom=2pt, left=3pt, right=3pt,
  toptitle=1pt, bottomtitle=1pt
}

% ── Hilfs-Befehle ───────────────────────────────────────────────────────
\newcommand{\regelnr}[1]{%
  \begin{tikzpicture}[baseline=-0.5ex]
    \node[circle, fill=arduinoBlue, text=white,
          font=\bfseries\small, inner sep=3pt] {#1};
  \end{tikzpicture}\enspace
}

\newcommand{\sectiontitle}[2]{%
  \noindent
  {\fontsize{9}{10}\selectfont\bfseries\color{arduinoBlue}%
   \regelnr{#1}#2}\\[-2pt]
  {\color{arduinoBlue!40}\rule{\linewidth}{0.6pt}}\\[2pt]
}

% ── Header ──────────────────────────────────────────────────────────────
% 1. Logo
% 2. Überschrift
% 3. Beschreibung
% 4. Zusatztext
\newcommand{\cheatheader}[4]{%
  \noindent
  \begin{tikzpicture}
    \node[fill=headerBg, text=white, rounded corners=4pt,
          minimum width=\linewidth, minimum height=1.0cm,
          inner sep=8pt, font=\bfseries] (h)
      {\fontsize{16}{18}\selectfont {#2}\quad
       \normalfont\fontsize{10}{12}\selectfont —\quad
       {#3}};
    \node[anchor=west, xshift=8pt] at (h.west)
    		{\includegraphics[height=0.75cm]{#1}};
    \node[anchor=east, xshift=-8pt] at (h.east)
      {\color{white!80}\fontsize{8}{9}\selectfont
       {#4}};
  \end{tikzpicture}
}

% Dokumentenversion
\newcommand\version{V1.0}

% ════════════════════════════════════════════════════════════════════════
\begin{document}
\pagestyle{empty}

% Header einfügen => Definition in cheatsheet.tex
\cheatheader{../images/fablab-logo.png}{Arduino Anfängerkurs}{Regeln für sauberen Code}{}

% Kurze Zusammenfassung der wichtigen Punkte
\noindent
\begin{tcolorbox}[
  colback=arduinoBlue!12,
  colframe=arduinoBlue,
  coltitle=white,
  fonttitle=\bfseries\fontsize{8}{9}\selectfont,
  title={\faStar\enspace Kurzübersicht — alle Regeln auf einen Blick},
  arc=3pt,
  boxrule=1pt,
  top=3pt, bottom=3pt, left=6pt, right=6pt,
  toptitle=2pt, bottomtitle=2pt
]
  \fontsize{8.5}{10}\selectfont
  \begin{tabular*}{\linewidth}{@{}l@{\extracolsep{\fill}}l@{\extracolsep{\fill}}l@{\extracolsep{\fill}}l@{}}
    \faTag\ \textbf{Pins:}\ \texttt{const int PIN\_X = N;}
    &
    \faFont\ \textbf{Namen:}\ \texttt{wartezeit},\ \texttt{ledIstAn}
    &
    \faComment\ \textbf{Kommentare} vor jedem Block
    &
    \faTerminal\ \textbf{Serial:}\ \texttt{Serial.begin(9600)}
    \texttt{+} \texttt{monitor\_speed = 9600}
  \end{tabular*}
\end{tcolorbox}

\begin{multicols}{3}
\raggedcolumns

% --------
% SPALTE 1
% --------

\sectiontitle{1}{Pins immer benennen}

Keine nackten Zahlen im Code! Jeder Pin bekommt am Anfang einen \textbf{Namen als Variable}.

\begin{badbox}
\begin{lstlisting}
pinMode(13, OUTPUT);
digitalWrite(13, HIGH);
\end{lstlisting}
\end{badbox}

\begin{goodbox}
\begin{lstlisting}
int LED_PIN = 13;

pinMode(LED_PIN, OUTPUT);
digitalWrite(LED_PIN, HIGH);
\end{lstlisting}
\end{goodbox}

% --------
\sectiontitle{2}{Sprechende Variablennamen}

Variablennamen müssen erklären, was sie speichern.

\begin{badbox}
\begin{lstlisting}
int x = 500;
int h = 200;
bool b = false;
\end{lstlisting}
\end{badbox}

\begin{goodbox}
\begin{lstlisting}
int  wartezeit  = 500; // ms
int  helligkeit = 200; // 0-255
bool ledIstAn   = false;
\end{lstlisting}
\end{goodbox}

\columnbreak

% --------
% SPALTE 2
% --------

\sectiontitle{3}{Blöcke kommentieren}

Jeder Block bekommt einen Kommentar, der \textbf{kurz} erklärt \textbf{was} hier passiert.

\begin{badbox}
\begin{lstlisting}
void loop() {
  int zustand = digitalRead(PIN_TASTER);

  if (zustand == HIGH) {
    digitalWrite(PIN_LED, HIGH);
  } else {
    digitalWrite(PIN_LED, LOW);
  }
}
\end{lstlisting}
\end{badbox}

\begin{goodbox}
\begin{lstlisting}
void loop() {
  // Taster auslesen
  int zustand = digitalRead(PIN_TASTER);

  // LED bei gedrücktem Taster einschalten
  // LED ausschalten, wenn Taster losgelassen wird^
  if (zustand == HIGH) {
    digitalWrite(PIN_LED, HIGH);
  } else {
    digitalWrite(PIN_LED, LOW);
  }
}
\end{lstlisting}
\end{goodbox}

\begin{regelbox}{ruleAccent}{\faLightbulb\enspace Kommentarstile}
  \fontsize{8}{9.5}\selectfont
  \texttt{// Einzeilig — kurze Erklärung}\\
  \texttt{/* Mehrzeilig für\\
  \phantom{xx}längere Erklärungen */}\\
  \texttt{int i = 0; // Am Zeilenende}
\end{regelbox}

\columnbreak

% --------
% SPALTE 3
% --------

\sectiontitle{4}{Serial Monitor einrichten}

\lstinline|Serial.begin(9600)| gehört in \textbf{jedes}
\lstinline|setup()| — auch wenn du ihn gerade nicht aktiv
nutzt. So kannst du jederzeit mit \lstinline|Serial.println()|
ausgeben lassen, was dein Arduino gerade macht.

\begin{goodbox}
\begin{lstlisting}
int LED_PIN    = 13;
int TASTER_PIN =  7;

void setup() {
  // Serial Monitor starten
  Serial.begin(9600);

  // Pins konfigurieren
  pinMode(LED_PIN,    OUTPUT);
  pinMode(TASTER_PIN, INPUT);

  Serial.println("Gestartet!");
}
\end{lstlisting}
\end{goodbox}

\vspace{4pt}

\begin{regelbox}{ruleAccent}{\faFile*[regular]\enspace platformio.ini}
  \fontsize{8}{9.5}\selectfont
  Die Baudrate im Code und in der Konfiguration \textbf{müssen übereinstimmen}!\\[3pt]
\begin{lstlisting}[language={},
  backgroundcolor=\color{white},
  basicstyle=\ttfamily\fontsize{7.5}{9}\selectfont]
[env:uno]
platform       = atmelavr
board          = uno
framework      = arduino
monitor_speed  = 9600
\end{lstlisting}
\end{regelbox}

\end{multicols}

% Footer
\vfill
\noindent
{\color{arduinoBlue!40}\rule{\linewidth}{0.5pt}}\\[1pt]
{\fontsize{7}{8}\selectfont\color{gray}
 \textbf{Arduino Anfängerkurs} — Code Regeln
 \hfill
 \today\ -\ \version
}

\end{document}